% Figure from MATH 493 notes, section 52. Compile with the notes preamble.
\begin{tikzpicture}[font=\small, every node/.style={inner sep=3pt}]
  \node (G) at (0,2.6) {$G = P_3P_5 \cong \mathbb{Z}/3\mathbb{Z} \times \mathbb{Z}/5\mathbb{Z} \cong \mathbb{Z}/15\mathbb{Z}$};
  \node[figR] (P3) at (-2.0,1.3) {$P_3 \trianglelefteq G$};
  \node[figB] (P5) at (2.0,1.3) {$P_5 \trianglelefteq G$};
  \node (E) at (0,0) {$P_3 \cap P_5 = \{e\}$};
  \draw (G.south) -- node[left, font=\scriptsize, black!60] {index $5$} (P3);
  \draw (G.south) -- node[right, font=\scriptsize, black!60] {index $3$} (P5);
  \draw (P3) -- node[left, font=\scriptsize, black!60] {$3$} (E);
  \draw (P5) -- node[right, font=\scriptsize, black!60] {$5$} (E);
  \node[font=\scriptsize, align=left, anchor=west] at (3.4,1.3) {$n_3 \in \{1,3,5,15\}$, $\equiv 1 \bmod 3$: $n_3 = 1$\\ $n_5 \in \{1,3,5,15\}$, $\equiv 1 \bmod 5$: $n_5 = 1$};
\end{tikzpicture}
